% Options for packages loaded elsewhere
\PassOptionsToPackage{unicode}{hyperref}
\PassOptionsToPackage{hyphens}{url}
\documentclass[
]{article}
\usepackage{xcolor}
\usepackage{amsmath,amssymb}
\setcounter{secnumdepth}{-\maxdimen} % remove section numbering
\usepackage{iftex}
\ifPDFTeX
  \usepackage[T1]{fontenc}
  \usepackage[utf8]{inputenc}
  \usepackage{textcomp} % provide euro and other symbols
\else % if luatex or xetex
  \usepackage{unicode-math} % this also loads fontspec
  \defaultfontfeatures{Scale=MatchLowercase}
  \defaultfontfeatures[\rmfamily]{Ligatures=TeX,Scale=1}
\fi
\usepackage{lmodern}
\ifPDFTeX\else
  % xetex/luatex font selection
\fi
% Use upquote if available, for straight quotes in verbatim environments
\IfFileExists{upquote.sty}{\usepackage{upquote}}{}
\IfFileExists{microtype.sty}{% use microtype if available
  \usepackage[]{microtype}
  \UseMicrotypeSet[protrusion]{basicmath} % disable protrusion for tt fonts
}{}
\makeatletter
\@ifundefined{KOMAClassName}{% if non-KOMA class
  \IfFileExists{parskip.sty}{%
    \usepackage{parskip}
  }{% else
    \setlength{\parindent}{0pt}
    \setlength{\parskip}{6pt plus 2pt minus 1pt}}
}{% if KOMA class
  \KOMAoptions{parskip=half}}
\makeatother
\usepackage{longtable,booktabs,array}
\usepackage{calc} % for calculating minipage widths
% Correct order of tables after \paragraph or \subparagraph
\usepackage{etoolbox}
\makeatletter
\patchcmd\longtable{\par}{\if@noskipsec\mbox{}\fi\par}{}{}
\makeatother
% Allow footnotes in longtable head/foot
\IfFileExists{footnotehyper.sty}{\usepackage{footnotehyper}}{\usepackage{footnote}}
\makesavenoteenv{longtable}
\setlength{\emergencystretch}{3em} % prevent overfull lines
\providecommand{\tightlist}{%
  \setlength{\itemsep}{0pt}\setlength{\parskip}{0pt}}
\usepackage{bookmark}
\IfFileExists{xurl.sty}{\usepackage{xurl}}{} % add URL line breaks if available
\urlstyle{same}
\hypersetup{
  hidelinks,
  pdfcreator={LaTeX via pandoc}}

\author{}
\date{}

\begin{document}

\section{\texorpdfstring{The Free \(\ell^1\) Seminorm on Banach Presheaf
Coboundaries}{The Free \textbackslash ell\^{}1 Seminorm on Banach Presheaf Coboundaries}}\label{the-free-ell1-seminorm-on-banach-presheaf-coboundaries}

\textbf{Author:} Jeremy H. Carroll \textbf{Date:} March 2026
\textbf{Version:} v012 (Pre-print)

\begin{center}\rule{0.5\linewidth}{0.5pt}\end{center}

\subsection{Abstract}\label{abstract}

We classify seminorms on the 1-cochain space of a Banach presheaf over a
finite poset satisfying additivity, faithfulness, and functoriality with
respect to bounded linear maps (contractive monotonicity). Specifically,
let \(P\) be a finite partially ordered set and \(F\) a Banach presheaf
over \(P\). The \textbf{1-cochain space} is
\(C^1(P, F) = \bigoplus_{e \in E_{\text{cov}}} F(a_e)\), where the
direct sum is taken in the \(\ell^1\) sense. The space \(C^1(P, F)\) is
the free additive Banach object generated by the edges. We prove that
the \(\ell^1\) coboundary norm \(I(s) = \sum_e \|\delta_e(s)\|\) is, up
to a positive scalar, the \textbf{unique edge-additive seminorm} on
\(C^1(P, F)\) satisfying these three properties. Adding isomorphism
invariance forces the scalar to be uniform across edges. No additional
choice remains: the \(\ell^1\) norm is uniquely determined by the
axioms, which encode compatibility with the underlying Banach geometry.

In short: \emph{any additive, faithful diagnostic on Banach presheaf
coboundaries that is functorial under arbitrary bounded linear maps must
be a positive scalar multiple of the \(\ell^1\) norm.}

\textbf{Conventions.} Norms are denoted \(\lVert \cdot \rVert\) with
appropriate subscripts; scalar absolute values are denoted \(|\cdot|\).

\begin{center}\rule{0.5\linewidth}{0.5pt}\end{center}

\subsection{1. Setup}\label{setup}

\subsubsection{1.1 Banach Presheaves over Finite
Posets}\label{banach-presheaves-over-finite-posets}

Let \((P, \leq)\) be a finite partially ordered set. The
\textbf{covering relations} are the edges of the Hasse diagram:
\[E_{\text{cov}}(P) = \{e = (a, b) : a < b,\; \nexists\, c \text{ with } a < c < b\}.\]

A \textbf{Banach presheaf} \(F\) over \(P\) assigns: - To each
\(x \in P\), a Banach space \((F(x), \|\cdot\|_{F(x)})\). - To each
\(e = (a, b) \in E_{\text{cov}}\), a bounded linear \textbf{restriction
map} \(r_{b,a}: F(b) \to F(a)\) with \(\|r_{b,a}\|_{\text{op}} \leq 1\).

A \textbf{section} of \(F\) is a family \(s = (s(x))_{x \in P}\) with
\(s(x) \in F(x)\).

\subsubsection{1.2 The 1-Cochain Space}\label{the-1-cochain-space}

For each edge \(e = (a, b)\), the \textbf{coboundary defect} is:
\[\delta_e(s) = r_{b,a}(s(b)) - s(a) \in F(a).\]

The \textbf{1-cochain space} is the Banach direct sum:
\[C^1(P, F) = \bigoplus_{e \in E_{\text{cov}}}^{\ell^1} F(a_e)\] with
norm \(\|(\delta_e)_e\|_1 = \sum_e \|\delta_e\|_{F(a_e)}\). This is
isometrically isomorphic to the space of coboundary vectors. As
discussed in Section 6.2, \(C^1(P, F)\) acts as the free additive Banach
object generated by the edges.

The \(\ell^1\) coboundary norm of a section \(s\) is:
\[I(s) = \sum_{e \in E_{\text{cov}}} \|\delta_e(s)\|_{F(a_e)} = \|(\delta_e(s))_e\|_1.\]

\subsubsection{1.3 Seminorms on
1-Cochains}\label{seminorms-on-1-cochains}

\begin{quote}
\textbf{Definition 1.1.} A \textbf{seminorm on 1-cochains} is a function
\(\nu: C^1(P, F) \to \mathbb{R}_{\geq 0}\) satisfying
\(\nu(cx) = |c|\nu(x)\) and \(\nu(x+y) \leq \nu(x) + \nu(y)\) for all
\(x,y \in C^1(P, F)\) and scalars \(c\).
\end{quote}

We consider four properties:

\textbf{(P1) Edge-wise additivity.} \(\nu\) decomposes as
\[\nu((\delta_e)_e) = \sum_{e \in E_{\text{cov}}} \nu_e(\delta_e)\]
where each \(\nu_e: F(a_e) \to \mathbb{R}_{\geq 0}\) is itself a
seminorm on \(F(a_e)\).

\textbf{(P2) Faithfulness.} For each \(e\):
\(\nu_e(\delta) = 0 \iff \delta = 0\).

\textbf{(P3) Functoriality under bounded linear maps (Contractive
Monotonicity).} The function \(\nu_e\) is compatible with the underlying
Banach geometry. Specifically, for every bounded linear endomorphism
\(\varphi: F(a_e) \to F(a_e)\) on the stalk space:
\[\nu_e(\varphi(\delta)) \leq \|\varphi\|_{\text{op}} \cdot \nu_e(\delta).\]
\emph{(This is a strong condition: it explicitly encodes functoriality
with respect to bounded linear endomorphisms of the individual stalk
spaces. Combined with Lemma 2.1, this forces the seminorm to depend only
on the Banach norm.)}

\textbf{(P4) Isomorphism invariance.} For any poset isomorphism
\(\pi: \tilde{P} \to P\):
\[\nu_{\tilde{P}}(\pi^*(\delta_e)_e) = \nu_P((\delta_e)_e).\]

\subsubsection{1.4 Relation to the Combinatorial
Classification}\label{relation-to-the-combinatorial-classification}

The present classification may be viewed as a functional-analytic
extension of the finite-dimensional result established in
``Coordinate-Wise Additivity and the \(\ell^1\) Norm on Finite Graph
Cochains'' {[}4{]}. In that setting, \(\ell^1\) arises from local
additivity over disjoint coordinate supports and edge separability
alone, with no Banach space structure required. Here, we replace
coordinate additivity with structural compatibility in the form of
functoriality under bounded linear maps on Banach spaces (P3), showing
that the same rigidity persists in a substantially more general setting.
The combinatorial result {[}4{]} thus provides the minimal case; the
present paper demonstrates that the \(\ell^1\) conclusion is not an
artifact of finite coordinates, but survives the passage to
infinite-dimensional Banach stalks.

\begin{center}\rule{0.5\linewidth}{0.5pt}\end{center}

\subsection{2. The Classification
Theorem}\label{the-classification-theorem}

We first isolate two structural lemmas.

\begin{quote}
\textbf{Lemma 2.0 (1-Dimensional Compatibility).} Let \(\nu_e\) satisfy
(P3). Then for any 1-dimensional Banach subspace
\(\text{span}(\delta) \subset V\), the restriction
\(\nu_e|_{\text{span}(\delta)}\) is well-defined and satisfies
(P1)--(P3) with respect to the induced Banach norm on
\(\text{span}(\delta) \cong \mathbb{R}\). Specifically: any isometric
embedding \(\iota: \mathbb{R} \hookrightarrow V\) induces a seminorm
\(\nu_e^{\mathbb{R}} = \nu_e \circ \iota\) on \(\mathbb{R}\) satisfying
(P2)--(P3), and any contractive surjection \(f: V \to \mathbb{R}\)
induces
\(\nu_e^{\mathbb{R}}(f(\delta)) \leq \|f\|_{\text{op}} \cdot \nu_e(\delta)\).

\emph{Proof.} Direct application of (P3) to \(\iota\) and \(f\)
respectively. Both are bounded linear maps, so (P3) applies without
additional hypothesis. \(\square\)
\end{quote}

\begin{quote}
\textbf{Lemma 2.1 (Direction Independence).} If \(\nu_e\) satisfies
(P3), then for any \(v, w \in F(a_e)\) with \(\|v\| = \|w\| \neq 0\), we
have \(\nu_e(v) = \nu_e(w)\).

\emph{Proof.} By the Hahn--Banach theorem, there exists
\(f_v \in F(a_e)^*\) with \(\|f_v\|_{\text{op}} = 1\) and
\(f_v(v) = \|v\|\). Define \(\Phi: F(a_e) \to F(a_e)\) by
\(\Phi(x) = f_v(x) \frac{w}{\|w\|}\). Because \(\Phi\) is a bounded
linear endomorphism of \(F(a_e)\), (P3) applies directly here. Then
\(\Phi(v) = w\) and \(\|\Phi\|_{\text{op}} = 1\), so
\(\nu_e(w) \leq \nu_e(v)\) by (P3). By symmetry, constructing
\(\Psi(x) = f_w(x) \frac{v}{\|v\|}\) gives \(\nu_e(v) \leq \nu_e(w)\),
forcing equality. \(\square\)
\end{quote}

\begin{quote}
\textbf{Theorem 2.2 (Free \(\ell^1\) Seminorm Classification).} Under
functoriality with respect to bounded linear maps, let \(\nu\) be a
seminorm on \(C^1(P, F)\) satisfying (P1), (P2), and (P3). Then for each
edge \(e\), there exists \(w_e > 0\) such that
\[\nu_e(\delta) = w_e \cdot \|\delta\|_{F(a_e)} \quad \text{for all } \delta \in F(a_e).\]

If additionally (P4) holds, then \(w_e = \alpha\) is independent of
\(e\), and
\[\nu((\delta_e)_e) = \alpha \sum_{e} \|\delta_e\|_{F(a_e)} = \alpha \cdot I(s).\]
\end{quote}

In other words: the \(\ell^1\) coboundary norm is the unique (up to
positive scalar) edge-additive seminorm satisfying all four properties.

\begin{center}\rule{0.5\linewidth}{0.5pt}\end{center}

\subsection{3. Proof}\label{proof}

The proof proceeds in three stages.

\subsubsection{Stage 1: Direction
Independence}\label{stage-1-direction-independence}

This is Lemma 2.1: for any two vectors \(v, w \in F(a_e)\) with
\(\|v\| = \|w\| \neq 0\), we have \(\nu_e(v) = \nu_e(w)\).

\subsubsection{\texorpdfstring{Stage 2: Each \(\nu_e\) is proportional
to
\(\|\cdot\|\).}{Stage 2: Each \textbackslash nu\_e is proportional to \textbackslash\textbar\textbackslash cdot\textbackslash\textbar.}}\label{stage-2-each-nu_e-is-proportional-to-cdot.}

Fix an edge \(e\) and let \(V = F(a_e)\). By Lemma 2.1, all unit vectors
in \(V\) have the same \(\nu_e\)-value: there exists \(w_e \geq 0\) such
that \(\nu_e(v) = w_e\) for all \(v\) with \(\|v\| = 1\). By (P2),
\(w_e > 0\).

For any \(\delta \in V\) with \(\delta \neq 0\), the absolute
homogeneity of \(\nu_e\) gives:
\[\nu_e(\delta) = \|\delta\| \cdot \nu_e\!\left(\frac{\delta}{\|\delta\|}\right) = \|\delta\| \cdot w_e = w_e \|\delta\|.\]

Since \(\nu_e(0) = 0\) by the seminorm property, we conclude
\(\nu_e(\delta) = w_e \|\delta\|\) for all \(\delta \in V\). \(\square\)
(Stage 2)

\subsubsection{Stage 3: Isomorphism invariance forces uniform
weights.}\label{stage-3-isomorphism-invariance-forces-uniform-weights.}

Because of (P1), the seminorm \textbf{splits as a direct sum}, so each
component can be evaluated in isolation. Let \(e = (a, b)\) be any
covering edge of a finite poset \(P\). Define the single-edge subposet
\(P_e = \{a < b\}\). Each summand \(\nu_e(\delta_e)\) depends only on
the coboundary defect \(\delta_e(s) = r_{b,a}(s(b)) - s(a)\), which is
computed entirely from the restricted presheaf data
\((F(a), F(b), r_{b,a})\). The value \(\nu_e(\delta_e)\) is therefore
intrinsically independent of the ambient poset \(P\) in which \(e\) is
embedded: it is determined by the single-edge presheaf \((P_e, F_e)\).

By Stages 1--2, \(\nu_e(\delta_e) = w_e \|\delta_e\|\) for some
\(w_e > 0\). Since \(w_e\) depends only on the isomorphism class of the
single-edge presheaf \((P_e, F_e)\), it remains to show all such
presheaves yield the same weight.

Let \(e_1 = (a_1, b_1)\) and \(e_2 = (a_2, b_2)\) be any two covering
edges with the same stalk spaces \(F(a_{e_1}) \cong F(a_{e_2})\) and
\(F(b_{e_1}) \cong F(b_{e_2})\), and compatible restriction maps under
the induced isomorphism. The single-edge posets
\(P_{e_1} = \{a_1 < b_1\}\) and \(P_{e_2} = \{a_2 < b_2\}\) admit a
unique order-preserving bijection \(\pi: P_{e_1} \to P_{e_2}\). By (P4),
\(\nu_{P_{e_1}}(\pi^*(\delta_{e_2})) = \nu_{P_{e_2}}((\delta_{e_2}))\),
which gives \(w_{e_1} = w_{e_2}\).

Thus all covering edges share the same weight \(w\). Setting
\(\alpha = w\): \(\nu = \alpha \cdot I\). \(\square\) (Theorem 2.2)

\begin{center}\rule{0.5\linewidth}{0.5pt}\end{center}

\subsection{4. Necessity of Each
Property}\label{necessity-of-each-property}

\begin{longtable}[]{@{}
  >{\raggedright\arraybackslash}p{(\linewidth - 4\tabcolsep) * \real{0.3333}}
  >{\raggedright\arraybackslash}p{(\linewidth - 4\tabcolsep) * \real{0.3333}}
  >{\raggedright\arraybackslash}p{(\linewidth - 4\tabcolsep) * \real{0.3333}}@{}}
\toprule\noalign{}
\begin{minipage}[b]{\linewidth}\raggedright
Dropped property
\end{minipage} & \begin{minipage}[b]{\linewidth}\raggedright
What happens
\end{minipage} & \begin{minipage}[b]{\linewidth}\raggedright
Example
\end{minipage} \\
\midrule\noalign{}
\endhead
\bottomrule\noalign{}
\endlastfoot
(P1) Additivity & \(\ell^\infty\)-type norms admitted &
\(\nu = \max_e \|\delta_e\|\) \\
(P2) Faithfulness & Trivial seminorm admitted & \(\nu \equiv 0\) \\
(P3) Functoriality & Direction-dependent seminorms admitted & On
\(F(a_e) = \mathbb{R}^2\):
\(\nu_e(x_1, x_2) = \lvert x_1\rvert + 2\lvert x_2\rvert\); satisfies
(P1), (P2), (P4), but violates direction independence (Lemma 2.1), hence
is not proportional to \(\|\cdot\|\) \\
(P4) Invariance & Non-uniform weights admitted &
\(\nu = \sum_e w_e \|\delta_e\|\), \(w_e\) edge-dependent \\
\end{longtable}

Each property is independently necessary. Together they are sufficient.

\begin{center}\rule{0.5\linewidth}{0.5pt}\end{center}

\subsection{5. Structural Uniqueness}\label{structural-uniqueness}

Rather than defining artificial categories to encapsulate simple scaling
relations, we formulate the universal structure directly:

\begin{quote}
\textbf{Proposition 5.1 (Normalized Uniqueness).} The \(\ell^1\)
coboundary norm is the unique normalized additive diagnostic satisfying
properties (P1)--(P4). Any evaluation compatible with these constraints
is identically \(\alpha \cdot I(s)\) for some positive scalar
\(\alpha\).
\end{quote}

\begin{quote}
\textbf{Corollary 5.2 (Commutative Factorization).} For any seminorm
\(\nu\) satisfying (P1)--(P4) and any bounded linear map \(\varphi\)
acting stalkwise, the evaluation functional factors predictably:

\[\nu(\varphi_*(\delta)) \leq \|\varphi\|_{\text{op}} \cdot \nu(\delta),\]

and the \(\ell^1\) coboundary norm \(I\) is the unique (up to scalar)
edge-additive framework for which this structural compatibility with
bounded linear maps identically holds.
\end{quote}

\begin{center}\rule{0.5\linewidth}{0.5pt}\end{center}

\subsection{6. Discussion}\label{discussion}

\subsubsection{6.1 What This Paper Proves}\label{what-this-paper-proves}

This paper isolates a purely structural constraint on how coboundary
defects in Banach presheaf systems can be measured: once one requires
edge-wise additivity, faithfulness, and compatibility with the ambient
Banach geometry through functoriality under bounded linear maps, the
resulting seminorm is forced to be \(\ell^1\) up to a global scalar. The
classification does not depend on any dynamical assumptions,
probabilistic structure, or physical interpretation; it is a consequence
of the algebraic and geometric properties encoded in (P1)--(P4). In
particular, the result shows that the freedom to choose a defect metric
on the 1-cochain space is more limited than it may appear ---
alternative geometries such as \(\ell^2\) or \(\ell^\infty\) are not
compatible under these constraints. The significance of the theorem is
therefore not the introduction of a new norm, but the identification of
conditions under which no alternative exists.

\subsubsection{6.2 Relation to Other Work}\label{relation-to-other-work}

The Hahn--Banach factoring argument used in Stage 1 is analogous to
classical norm classification techniques in Banach space theory, where
linear functionals are used to reduce geometric constraints to
one-dimensional analysis (see Megginson 1998 {[}2{]} for Banach
geometric tools, and Pestov 2006 {[}1{]} for related uniqueness results
for norms). The key structural observation is that functoriality under
all bounded linear endomorphisms, combined with Hahn--Banach, forces
direction independence. This mechanism parallels the uniqueness proof
for the operator norm.

The result is a presheaf version of a classical principle in Banach
space theory: the characterization of the \(\ell^1\) norm as the unique
additive cross-norm on a free Banach sum. The 1-cochain space
\(C^1(P, F) = \bigoplus_e F(a_e)\) is the \emph{free additive Banach
object generated by the edges} of the poset. When a measurement on such
a free object must satisfy piecewise additivity and strict compatibility
with the underlying norms on each component, \(\ell^1\) is the only
geometry that survives. The \(\ell^2\) norm fails because
\(\|(v, w)\|_2 \neq \|v\| + \|w\|\); the \(\ell^\infty\) norm fails
because \(\max(\|v\|, \|w\|)\) discards summation. This positions the
\(\ell^1\) coboundary norm as the canonical additive cross-norm on the
free Banach object \(C^1(P,F)\) under the stated constraints.

We note that the coboundary defects
\(\delta_e(s) = r_{b,a}(s(b)) - s(a)\) measure the failure of local data
to glue consistently across the covering relations of the poset. The
distinction in our framework is that the choice of how to measure these
defects is not arbitrary: it is uniquely forced to be the \(\ell^1\)
coboundary norm by the structural axioms (P1)--(P4).

\subsubsection{6.3 Strategic Place in the Overall
Program}\label{strategic-place-in-the-overall-program}

This paper sits as Step 2 within a four-paper sequence establishing that
additive, structure-preserving measurements of local defects are
uniquely forced to be \(\ell^1\) across progressively richer structural
settings, and that projection dynamics generically produce nonzero
instances of such defects. The sequence progresses: 1.
\textbf{Coordinate-Wise Additivity} {[}4{]}: Local additivity over
disjoint coordinate supports on finite graph cochains already forces
\(\ell^1\) geometry. This is the purely combinatorial core. 2.
\textbf{Free \(\ell^1\) Seminorm Classification} (this paper): The same
\(\ell^1\) rigidity persists when the cochain spaces are replaced by
Banach presheaves and coordinate additivity is replaced by structural
compatibility (functoriality under bounded linear maps). 3.
\textbf{Single Obstruction Theory} {[}3{]}: Dynamical retractions
produce intrinsic coboundary defects measured by the forced \(\ell^1\)
diagnostic, and these defects persist with non-vanishing density over
finite-time intervals protected by Lieb--Robinson bounds. 4.
\textbf{Hodge Decomposition} {[}5{]}: The forced \(\ell^1\) geometry
induces a canonical quotient norm on \(H^1(G)\), providing the unique
gauge-invariant measure of irreducible obstruction.

\subsubsection{6.4 Connection to Dynamical Obstruction
Theory}\label{connection-to-dynamical-obstruction-theory}

The classification result established here is purely structural: it
identifies the unique form of any additive, functorial diagnostic on
presheaf coboundaries. Its significance becomes more apparent when
combined with dynamical settings in which such coboundary defects are
generated.

In the companion work ``Projection Obstruction Theory'' {[}3{]},
projection of linear many-body dynamics onto a nonlinear manifold is
shown to produce intrinsic coboundary defects that cannot be removed by
linearization and that persist with non-vanishing density over finite
time intervals. The present result implies that any local, additive
measurement of such defects must coincide with the \(\ell^1\) coboundary
norm (up to scaling). Thus, the \(\ell^1\) diagnostic appearing in that
analysis is not a modeling choice, but the only one compatible with the
structural constraints imposed here.

\emph{Observer-theoretic grounding.} Properties (P1)--(P4) admit a
natural operational interpretation: they correspond to the structural
conditions under which a local agent can consistently detect and
aggregate coboundary defects. (P1) requires local evaluability; (P2)
requires nontriviality; (P3) requires that local repairs do not increase
total measured defect; (P4) requires indifference to relabeling. Under
this reading, Theorem 2.2 states that any defect diagnostic usable by a
consistent local observer on Banach presheaf coboundaries must be
proportional to the \(\ell^1\) coboundary norm. This interpretation is
developed in {[}3, Section 1.2{]} and parallels the combinatorial
version established in {[}4, Section 3.3{]}.

\begin{center}\rule{0.5\linewidth}{0.5pt}\end{center}

\subsection{Acknowledgments}\label{acknowledgments}

No external funding. No conflicts of interest.

\begin{center}\rule{0.5\linewidth}{0.5pt}\end{center}

\subsection{References}\label{references}

{[}1{]} V. Pestov, ``Dynamics of infinite-dimensional groups and
Ramsey-type phenomena,'' \emph{Publicacoes Matematicas do IMPA}, 2006.
{[}2{]} R. E. Megginson, \emph{An Introduction to Banach Space Theory},
Springer, 1998. {[}3{]} J. H. Carroll, ``Projection Obstruction Theory:
Retraction Nonlinearity, \(\ell^1\) Rigidity, and Density Scaling,''
Zenodo, 2026. https://doi.org/10.5281/zenodo.18896776 {[}4{]} J. H.
Carroll, ``Coordinate-Wise Additivity and the \(\ell^1\) Norm on Finite
Graph Cochains,'' Zenodo, 2026. https://doi.org/10.5281/zenodo.18945342
{[}5{]} J. H. Carroll, ``Hodge Structure and Gauge Equivalence of
\(\ell^1\) Defect Fields,'' Zenodo, 2026.

\end{document}
